\section{Discussion}
\subsection{Derived FOPDT transfer functions.}
\subsubsection{Matlab fitting}
The \textbf{tfest} function was used in Matlab to fit the data to a transfer function. The following steps were followed to achieve this:
\\
\\
\textbf{Data Preparation:}
\\
The recorded data was sliced to include a sufficient amount of zero inputs, this is required by the tfest function to properly fit the model. After which the data was broken into a voltage vector (the PWM input), and two temperature vectors 
(temperatures of each of the transistors). The ambient temperature was subtracted from  the temperature vectors and the data was packaged into an iddata object.
\\
\\
\textbf{Model Fitting:}
\\
Three parameters were passed to the tfest function, namely: the iddata object, the number of poles the system has, and the number of zeros the system has.
The number of poles a system has is determined by the number of energy storing elements, in the case of this system the transistors are storing energy in the form
of heat, thus the system has two poles. It is common practise to take the number of zeros of an unknown system as: Zeros $=$ Poles $- 1$.


